\documentclass[a4paper]{article}
\usepackage{ctex}
\ctexset{proofname=\heiti{证明}}
\usepackage{amsmath,amssymb,amsthm}
\usepackage{mathtools}
\usepackage{bm}
\usepackage{enumitem}
\usepackage{booktabs}
\usepackage{array}
\usepackage{graphicx}

\IfFileExists{iidef.sty}{
    \usepackage[thehwcnt = 1]{iidef}
    \thecourseinstitute{清华大学}
    \thecoursename{人工智能原理}
    \theterm{2026年春季学期}
    \hwname{作业}
    \slname{\heiti{解}}
}{
    \makeatletter
    \newcommand{\thecourseinstitute}[1]{\def\@courseinstitute{#1}}
    \newcommand{\thecoursename}[1]{\def\@coursename{#1}}
    \newcommand{\theterm}[1]{\def\@term{#1}}
    \newcommand{\hwname}[1]{\def\@hwname{#1}}
    \newcommand{\slname}[1]{\def\@slname{#1}}
    \thecourseinstitute{清华大学}
    \thecoursename{人工智能原理}
    \theterm{2026年春季学期}
    \hwname{作业}
    \slname{\heiti{解}}
    \newcommand{\courseheader}{
        \begin{center}
            {\Large \heiti \@courseinstitute \quad \@coursename}\\[0.5em]
            {\large \@term \quad \@hwname 1 \quad \@slname}
        \end{center}
        \vspace{1em}
    }
    \newcommand{\name}[1]{\noindent #1 \par\vspace{1em}}
    \newenvironment{solution}{\par\noindent\textbf{解：}\par}{\par}
    \makeatother
}

\begin{document}
\courseheader
\name{姓名：杜一郎\qquad 学号：2023011618}

\begin{enumerate}
\setlength{\itemsep}{3\parskip}

\item 搜索问题形式化

\begin{solution}
\textbf{(a) 四枚硬币翻转}

状态可以写成四元组
\[
s=(x_1,x_2,x_3,x_4),\qquad x_i\in\{\text{正},\text{反}\}.
\]
初始状态为
\[
s_0=(\text{正},\text{反},\text{正},\text{反}),
\]
目标状态为
\[
s_g=(\text{正},\text{正},\text{正},\text{正}).
\]
每次动作选择连续的三枚硬币并翻转，可选动作是
\[
\mathrm{flip}(1,2,3),\qquad \mathrm{flip}(2,3,4).
\]
转移模型为：被选中的硬币正反互换，未被选中的硬币保持不变。每次翻转代价为 1，总路径代价就是翻转次数。

\textbf{(b) 网格机器人}

状态可设为机器人所在格子坐标
\[
s=(i,j).
\]
若题目只要求找一条从起点到终点的路径，状态中不需要记录历史路径。初始状态为
\[
S=(1,1),
\]
目标状态为
\[
T=(N,N).
\]
动作集合为
\[
\{\text{上},\text{下},\text{左},\text{右}\}.
\]
若移动后的格子仍在边界内且不是障碍格 \(O\)，则机器人移动到该格；否则该动作非法。每移动一格代价为 1，因此总代价为路径长度。
\end{solution}

\item 城市路网搜索

\begin{solution}
根据图中边权，主要边为
\[
\begin{gathered}
S\!-\!A=4,\ S\!-\!B=3,\ S\!-\!C=5,\ A\!-\!B=2,\ B\!-\!C=2,\\
A\!-\!D=5,\ B\!-\!D=4,\ B\!-\!E=3,\ C\!-\!F=4,\\
D\!-\!E=2,\ D\!-\!H=3,\ E\!-\!H=4,\ E\!-\!K=4,\ E\!-\!F=3,\\
F\!-\!K=5,\ H\!-\!T=5,\ K\!-\!T=4.
\end{gathered}
\]

\textbf{(a) BFS}

按字母序扩展相邻节点，BFS 的扩展顺序为
\[
S,A,B,C,D,E,F,H,K,T.
\]
对应搜索树的父节点如下：
\[
\begin{array}{c|ccccccccc}
\toprule
\text{节点} & A & B & C & D & E & F & H & K & T\\
\midrule
\text{父节点} & S & S & S & A & B & C & D & E & H\\
\bottomrule
\end{array}
\]
因此 BFS 找到的路径为
\[
S\rightarrow A\rightarrow D\rightarrow H\rightarrow T.
\]
BFS 只按层数搜索，不考虑边权，所以这条路径不一定是总代价最小的路径。

\textbf{(b) UCS}

UCS 每次扩展累计代价 \(g(n)\) 最小的节点。扩展顺序为
\[
\begin{array}{c|c|c}
\toprule
\text{顺序} & \text{节点} & g(n)\\
\midrule
1&S&0\\
2&B&3\\
3&A&4\\
4&C&5\\
5&E&6\\
6&D&7\\
7&F&9\\
8&H&10\\
9&K&10\\
10&T&14\\
\bottomrule
\end{array}
\]
得到最短路径
\[
S\rightarrow B\rightarrow E\rightarrow K\rightarrow T,
\]
总代价为
\[
3+3+4+4=14.
\]

\textbf{(c) A*}

使用题目给出的启发函数，令
\[
f(n)=g(n)+h(n).
\]
按 \(f\) 最小扩展，平局按字母序处理。主要过程如下：
\[
\begin{array}{c|c|c|l}
\toprule
\text{步} & \text{扩展节点} & g & \text{Open List 主要节点 }(g,f)\\
\midrule
1&S&0&B(3,12), A(4,15), C(5,16)\\
2&B&3&E(6,12), D(7,14), A(4,15), C(5,16)\\
3&E&6&K(10,13), D(7,14), H(10,14), A(4,15), C(5,16), F(9,17)\\
4&K&10&D(7,14), H(10,14), T(14,14), A(4,15), C(5,16), F(9,17)\\
5&D&7&H(10,14), T(14,14), A(4,15), C(5,16), F(9,17)\\
6&H&10&T(14,14), A(4,15), C(5,16), F(9,17)\\
7&T&14&\text{结束}\\
\bottomrule
\end{array}
\]
最终路径同样为
\[
S\rightarrow B\rightarrow E\rightarrow K\rightarrow T,
\]
总代价为 \(14\)。
\end{solution}

\end{enumerate}
\end{document}
